\begin{tikzpicture}[
  >={Latex[length=3mm,width=2.2mm]},
  every node/.style={font=\small},
  loop/.style={->,line width=0.7pt,loopblue},
  mvec/.style={->,line width=1.3pt,mred}
]

% Sagoma superiore del materiale magnetizzato
\def\blobtop{%
  plot [smooth cycle, tension=0.9] coordinates {
    (-4.2,1.6) (-2.0,2.6) (0.4,2.4) (2.4,2.7) (4.4,2.0)
    (5.2,0.6) (4.0,-0.4) (2.8,0.2) (1.2,-0.2) (-0.6,0.1)
    (-2.4,-0.2) (-3.8,0.2)
  }
}

% Faccia inferiore leggermente traslata verso il basso (spessore t)
\begin{scope}[yshift=-0.6cm]
  \fill[matpinkdk] \blobtop;
\end{scope}

% Linea verticale sui bordi che collega superficie sup e inf (spessore)
\draw[matpinkdk,line width=0.4pt] (-4.2,1.6) -- ++(0,-0.6);
\draw[matpinkdk,line width=0.4pt] (5.2,0.6) -- ++(0,-0.6);

% Faccia superiore
\fill[matpink,draw=matpink!60!black,line width=0.4pt] \blobtop;

% Indicazione spessore t sulla destra
\draw[->,thin] (5.4,0.9) -- (5.4,0.3);
\draw[->,thin] (5.4,-0.9) -- (5.4,-0.3);
\node[right] at (5.45,0.0) {$t$};

% Griglia di piccoli loop di corrente (rappresentano le spire magnetiche)
\foreach \x/\y in {%
  -3.2/1.6, -2.2/1.9, -1.2/2.0, -0.2/1.9, 0.8/1.9, 1.8/2.0, 2.8/1.9, 3.6/1.7,
  -3.4/0.9, -2.4/1.1, -1.4/1.2, -0.4/1.2, 0.6/1.2, 1.6/1.2, 2.6/1.1, 3.4/0.9,
  -3.0/0.2, -2.0/0.4, -1.0/0.5, 0.0/0.5, 1.0/0.5, 2.0/0.5, 3.0/0.3%
}{
  \begin{scope}[shift={(\x,\y)}, scale=0.35]
    % piccolo rombo con freccia (loop di corrente)
    \draw[loop] (-0.5,0) -- (0,0.3) -- (0.5,0) -- (0,-0.3) -- cycle;
    \draw[loop] (-0.35,-0.05) -- (0.05,0.2);
  \end{scope}
}

% Etichette I sparse
\node[loopblue] at (-3.5,2.3) {$I$};
\node[loopblue] at (-0.1,2.4) {$I$};
\node[loopblue] at ( 1.6,0.7) {$I$};
\node[loopblue] at (-2.3,0.6) {$I$};
\node[loopblue] at ( 2.9,1.5) {$I$};

% Vettore magnetizzazione M (al centro, verso l'alto)
\draw[mvec] (3.0,0.8) -- (3.0,2.8) node[right] {$\mathbf{M}$};

\end{tikzpicture}